\section{The remaining equations}
\label{sec:targets}

We now state the algebraic target for the contour calculation.
The coefficient formula \eqref{foundation:core} can also be written
\begin{equation}\label{tail:uniform-core}
 \mathsf B_{a,b}=\delta^{-1}W_aW_{b-a}W_{-b}
                  +\tau^2\delta^{-1}\ind{a=b=0}.
\end{equation}
It immediately gives the linear reconstruction identities
\begin{equation}\label{tail:linear}
 \mathsf A_{a,b}=\mathsf A_{-b,a-b},\qquad
 \mathsf A_{a,0}=\mathsf A_{0,a}=\ind{a=0}-\tau^{-1},
 \qquad \sum_a\mathsf A_{a,0}=-\delta^{-1}.
\end{equation}
Define the residuals
\begin{align}
 Q(g,h)&=\sum_m\mathsf A_{m+g,h}\mathsf A_{h,m}
                  -\ind{g=0}+\delta^{-1}\ind{h=0},\label{tail:Q}\\
 C(g,h,k,l)&=\sum_m\mathsf A_{m,g+h}\mathsf A_{g,m+k}
                         \mathsf A_{h,m+l}
       -\mathsf A_{g+l,k}\mathsf A_{h+k,l}
       +\delta^{-1}\ind{g=h=0},\label{tail:C}\\
 \mathcal C_{g,h}(k,l)&=\tau^3C(g,h,k,l).\label{tail:scaledC}
\end{align}
Proposition~\ref{foundation:quadratics} says $Q=0$.
Parts~\ref{part:contours} and~\ref{part:completion} prove $C=0$;
Part~\ref{part:reconstruction} then derives the quartic equations.
The two normalizations $C$ and $\mathcal C$ express the same condition:
$C$ uses the reconstruction matrix $\mathsf A$, while $\mathcal C$
avoids denominators in calculations with $\mathsf B$.
